\documentclass[12pt]{article}
\usepackage[paperwidth=6in,paperheight=9in]{geometry}
\geometry{
  inner=15mm,
  outer=10mm,
  top=13mm,
  bottom=13mm,
  headheight=12pt,
  headsep=7mm,
  footskip=10mm,
  includehead=true,
  includefoot=true
}
\usepackage{fontspec}
\setmainfont{Libertinus Serif}
\usepackage{microtype}
\usepackage{setspace}
\usepackage{ragged2e}

\title{Chapter 15: The Defenders and Their Strategies}
\author{From \textit{Stolen Words}}
\date{}

\raggedbottom
\clubpenalty=150
\widowpenalty=150
\setlength{\parskip}{4pt plus 2pt minus 1pt}
\setlength{\parindent}{1em}

\begin{document}

\maketitle

\vspace{1cm}

{\centering\itshape When institutions say 'these are minor improvements,'\\they're asking you to trust their judgment\\over your own spiritual experience.\par}
\vspace{0.5cm}

\justifying

Maya Rodriguez knew her investigation would eventually lead here—to the defenders of the revision. After David Matthews revealed the publishing deception, she needed to understand how institutions justified what had been done.

She arranged a meeting with Dr. Richard Whitfield (Ritambhara Dasa), a senior BBT representative who had publicly defended the revisions for two decades. They met at the BBT offices in Los Angeles, a modern building filled with Sanskrit texts and photographs of Prabhupāda.

As Maya entered Whitfield's office, she noticed the man's genuine reverence—photos of him with various spiritual teachers, Sanskrit dictionaries worn from use, devotional texts in multiple languages. This wasn't a corporate executive; this was someone who had dedicated his life to spiritual service.

"Ms. Rodriguez," Dr. Whitfield greeted her formally, but Maya caught something in his eyes—perhaps a flicker of the same uncertainty she'd seen in her own mirror. "I understand you have questions about our editorial process."

It struck Maya that this man had probably asked himself the same questions she was asking him.

Maya opened her notebook, now thick with documentation. "I have evidence of systematic alteration affecting the overwhelming majority of verses. How do you justify this?"

What followed would be a masterclass in institutional defense mechanisms.


\vspace{0.5cm}
\textbf{The Defense}
\vspace{0.2cm}


The first fifteen minutes were predictable institutional deflection. Dr. Whitfield employed every standard defense:

"These are minor editorial improvements, not substantial changes."

Maya spread her statistical analysis across his desk. "Five hundred forty-one verses out of seven hundred. You call that minor?"

"We improved Sanskrit accuracy, scholarly apparatus, editorial professionalism—"

"You packaged technical improvements with theological revision," Maya interrupted. "Why couldn't you fix diacritical marks without changing 'Blessed Lord' to 'Supreme Personality of Godhead'? You could have created a 'Scholar's Edition,' clearly labeled. Instead, you replaced the original and hid the changes."

"The institution authorized these revisions. The GBC approved—"

Maya pulled out a photograph of Prabhupāda. "This man had spiritual realization. He chose specific words for specific reasons. Your committees had what—good English degrees? There's a difference between administrative competence and spiritual realization. You've confused the two."

Dr. Whitfield's jaw tightened. But Maya saw something shift in his eyes—not surrender, but recognition that she had done her homework.


\vspace{0.5cm}
\textbf{Defense Strategy 4: The "Prabhupāda Wanted Revisions" Defense}
\vspace{0.2cm}


Dr. Whitfield played his strongest card: "Prabhupāda wanted these changes but didn't have time to implement them."

Maya had been waiting for this claim. She pulled out a thick folder labeled "Class Transcripts."

"Let's examine your claims one by one," Maya said. "You say there were unpublished instructions. Can you show me one letter where Prabhupāda asked for 'Blessed Lord' to be changed to 'Supreme Personality of Godhead'?"

"The principle was established through his corrections to other works—"

"But we're talking about the Gita specifically. And these draft preferences—drafts are drafts, Dr. Whitfield. Publication represents the author's final decision, doesn't it?"

Whitfield hesitated. "In ideal circumstances, yes, but—"

"You also mentioned time constraints. He had five years from 1972 until his passing in 1977. Are you saying that wasn't enough time?"

"The situation was complex—"

"Then explain this: if he was such a perfectionist about these texts, why did he approve the 1972 edition for publication?"

Dr. Whitfield shifted through his papers. "We have manuscript evidence—"

"Stop," Maya said firmly. "Let me show you what I have."

She opened the class transcripts:

\textbf{"Five Years of Published Use: From 1972 to 1977, Prabhupāda used his published Bhagavad-gītā in hundreds of classes. Not once—not once—did he request the changes you made."}

Maya showed specific examples:

\textbf{December 16, 1968, Los Angeles}: "Listen to this transcript. A devotee reads verse 2.48 with 'steadfast in yoga' and 'evenness of mind.' Prabhupāda's response? He emphasizes these exact concepts. No correction."

\textbf{1974 Lecture}: "Here, verse 6.31 is read with 'worships Me.' Prabhupāda emphasizes: 'This is bhakti-yoga. Not impersonal meditation—direct worship of Krishna.' He's approving direct personal worship. You later redirected it to 'worshipful service of the Supersoul.'"

\textbf{March 1975, Māyāpur}: "Verse 2.30 read with 'eternal.' Prabhupāda repeats 'eternal' five times in his explanation. You removed it."

Dr. Whitfield was sweating now.

"If Prabhupāda wanted these changes," Maya pressed, "where are the letters requesting specific alterations? Where are the class corrections when verses were read? Where are the instructions to editors about improvements? Where are the meeting notes with revision requests?"

"He mentioned things privately—"

"Privately to whom? Where's the documentation? You've changed a published book based on undocumented private conversations?"

Dr. Whitfield pulled out his final argument: "Modern drafts reveal Prabhupāda's true theological intentions. We have manuscripts showing what he really wanted."

Maya had researched this claim thoroughly.

"Let me understand your position," Maya said. "Unpublished drafts override published books?"

She leaned forward. "You know better than Prabhupāda what he meant?"

Whitfield's fingers tightened on his papers.

"He would approve changes he never requested?"

"Secret drafts reveal secret intentions?"

"He wanted changes but never said so?"

A long silence. When Whitfield finally spoke, his voice had lost its earlier confidence. "The manuscripts show—"

\textbf{The Primary Source Contradiction}

Maya interrupted: "Dr. Whitfield, you're an educated man. In any field—history, literature, science—what's the primary source?"

"The original document, but—"

"The published work or the draft?"

Silence.

"When Prabhupāda published the Bhagavad-gītā in 1972, that was his final editorial decision. That's what he chose to give the world. You're saying unpublished drafts override published decisions?"

\textit{Draft Irrelevance and Selective Evidence}

Maya pulled out her research on the drafts: "You found isolated instances where Prabhupāda crossed out certain phrases in drafts. From this, you concluded systematic theological revision was authorized?"

"The pattern was clear—"

"The pattern? He used 'Blessed Lord' in the published book! He taught from it for five years! That's the pattern!"

She continued: "Every author has drafts with crossed-out words. The publication is what they decided to keep. You're cherry-picking draft evidence while ignoring five years of him using the published version."

\textit{Cultural Precedent Violation}

"Dr. Whitfield," Maya asked, "would you rewrite Shakespeare because you found a draft where he crossed out 'To be or not to be'?"

"That's different—"

"How? Would you 'improve' Beethoven's Ninth Symphony because you found a rejected draft?"

"Religious texts—"

"Are somehow less deserving of preservation? If anything, they deserve more protection, not less."

\textit{The Authentication Problem}

"Here's what I don't understand," Maya said. "If you believe your version is better, why not be honest about it? Call it 'Bhagavad-gītā As It Is: BBT Revised Edition.' Let people choose."

Dr. Whitfield's answer revealed everything: "That would confuse people."

"No," Maya replied. "It would inform them. And that's what you're afraid of."

Maya showed her final evidence: "When Prabhupāda wanted changes, look at his pattern:"

She read from his letters:
\begin{itemize}
\item 1970: "I am sending the necessary Sanskrit corrections"
\item 1971: "So when these corrections are made then you can print"
\item 1973: "That 'regulated' should be 'rejected'—please correct"
\end{itemize}

"Immediate. Specific. Clear. If he wanted systematic divine address changes, he had 1,825 days to request them. He didn't."

Dr. Whitfield was silent for a long moment. Maya watched him staring at the photograph of Prabhupāda on his desk—a picture of his spiritual master, the man whose words he had spent twenty years defending alterations to.

When he looked up, his institutional mask had slipped slightly.

"You know," he said quietly, "there are nights I lie awake wondering if we made the right choice. Twenty years ago, I believed completely that we were serving Prabhupāda by perfecting his work. Now..."

He paused, seeming to weigh something internally.

"But then I think about the thousands of people who've found spiritual life through our version. Are you asking me to tell them their spiritual development is invalid?"

Maya realized she was seeing the human cost of institutional positions—not just on readers, but on the defenders themselves.

"Ms. Rodriguez," he said softly, "you're very thorough. But you're missing the bigger picture."

"Which is?"

"We're not destroying Prabhupāda's work. We're fulfilling it. The revision represents his mature theological vision—what he would have wanted if he'd had time."

Maya recognized the core presumption at the heart of the entire controversy. But Whitfield wasn't finished defending his position.


\vspace{0.5cm}
\textbf{The Final Assault}
\vspace{0.2cm}


Dr. Whitfield tried one more approach, leaning forward with genuine conviction. "Ms. Rodriguez, you're looking at this through an idealized lens. Let me give you context."

He spoke of Prabhupāda's final years—illness, institutional pressures, thousands of pages of manuscripts. "Should we have abandoned that material? Or completed what he intended but couldn't finish?"

He showed her photographs of manuscript pages with Prabhupāda's handwriting. "These aren't drafts. These are late-stage proofs. Every writer refines their work. We preserved his later refinements."

"Did he ask you to replace the published edition with these manuscripts?" Maya asked quietly. "Did he say, 'After I die, use these instead of what I approved'?"

Whitfield hesitated.

"That's what I thought," Maya said. "You made that decision for him."

Dr. Whitfield pulled out his final card: "With all due respect, you're not initiated. You're an outside academic. Do you think you understand Prabhupāda's intentions better than his direct disciples?" He gestured to a photo showing himself as a young man with Prabhupāda. "I was there. I heard him speak about exact Sanskrit translation, about trusting editors to perfect his English. You've read transcripts. We lived it."

Maya recognized the move—an appeal to authority, the classic fallacy when evidence runs thin.

"Dr. Whitfield," she said quietly, "you're right. I'm not initiated. I didn't know Prabhupāda personally. But that's exactly why we need objective evidence."

She leaned forward. "If his direct disciples disagree about what he wanted—and they clearly do, given the lawsuits—then the only neutral arbiter is what he actually published and used for five years. Not what you remember him saying. Not what you think he meant. What he approved in print and taught from repeatedly."

"I can't interpret his spiritual realizations. But I can read. He chose 'Blessed Lord.' He chose 'steadfast in yoga.' He taught from those choices for five years. Never once did he request them changed."

Maya met his eyes. "Here's my question: if your justifications are so strong, why not publish both versions, side by side, and let readers choose? Why not include an introduction explaining your reasoning?"

"You genuinely believe you're serving his mission. I respect that. But what if serving his mission means preserving his words, not improving them?"

\textbf{The Return to Institutional Defense}

Dr. Whitfield's expression hardened. "I think we're done here, Ms. Rodriguez. You've clearly already decided what you think."

"I'm just asking you to look at the evidence—"

"Evidence?" He stood, gathering papers. "You have correlation, not causation. You have preferences, not proof. Different people respond to different styles of spiritual writing—that doesn't mean one translation is 'programming' anyone."

"But the systematic nature of the changes—"

"Demonstrates systematic editorial effort to improve fidelity to the Sanskrit. Which is exactly what we've said publicly for forty years." He moved toward the door. "I'm afraid I have another meeting."

Maya felt the conversation slipping away. "Dr. Whitfield, off the record—don't you think readers deserve to know both versions exist? Don't they deserve conscious choice?"

He paused, hand on the doorknob. For a moment, something shifted in his face—not confession, but perhaps recognition.

"Ms. Rodriguez, the BBT's position is documented in our published responses. Jayadvaita Swami has written extensively about the revision process. Our editorial philosophy is transparent: we believe the revised edition more faithfully represents Śrīla Prabhupāda's Sanskrit understanding. If you disagree, publish your findings. That's how scholarship works."

"Will you respond to my findings?"

"The BBT responds to scholarly criticism in scholarly forums. If you publish, we'll consider it." He opened the door. "Good day, Ms. Rodriguez."

After leaving Dr. Whitfield's office, Maya sat in her car, frustrated but not surprised. She'd gotten the official position—professional, defensive, giving nothing away. No dramatic confession. No smoking gun. Just institutional stonewalling dressed in scholarly language.

But his phrase haunted her: "If you disagree, publish your findings."

Perhaps that was the point. Perhaps the battle wasn't about convincing institutions to acknowledge what they'd done. Perhaps it was about giving readers the information to make their own choices.

She called Dr. Sarah Chen at Stanford.

"Sarah, I think I need to write this as a book, not a dissertation. The BBT won't engage. But readers will."

"What did Whitfield say?"

"Exactly what you'd expect. Nothing I can use. Everything defensible. That's the problem with institutions—they don't confess. They just... persist."

There was silence on the line.

"So what now?" Chen asked.

"Now I write it as a book. Not for the BBT—they won't engage. Not even for my dissertation committee. For readers. For people like my grandmother who deserve to know they have a choice."

"Do you have enough evidence?"

"I have nine months of documentation, temple ethnography, the neuroscience research, class transcripts, statistical analysis. I have everything except a confession—and apparently that's not how institutions work."

"Then write it," Chen said. "Make the case. Let readers decide."

As Maya drove home, she thought about the millions of readers worldwide, unknowingly choosing between two spiritual orientations based on which version they happened to purchase. No dramatic conspiracy. No villainous confession. Just systematic editorial changes implemented by sincere people who believed they were improving fidelity to Sanskrit—and in the process, transformed how readers encountered the divine.

The answer wouldn't come from institutions acknowledging what they'd done. It would come from readers making conscious, informed choices about their own spiritual development.

\end{document}
